% Inserir em Sobre o modelo de referência. A chave bibliográfica está em ito.bib.
% Números de arquitetura: metadados do torchvision 0.29.2.
\subsection{MobileNetV3-Large}

A MobileNetV3~\cite{howard2019mobilenetv3} é uma rede feita para rodar em
celulares e sistemas embarcados, onde o custo de processamento importa tanto
quanto a acurácia. Ela reaproveita as convoluções separáveis em profundidade
(\textit{depthwise separable}) da MobileNetV1 e os blocos residuais invertidos
da MobileNetV2. A isso foram somados módulos \textit{squeeze-and-excitation} e
a ativação \textit{hard-swish}, uma versão mais barata de calcular da
\textit{swish}. A ordem e o tamanho dos blocos foram definidos por busca
automática de arquitetura (NAS), e não à mão.

Usamos a versão \textit{Large} do Torchvision com os pesos
\texttt{IMAGENET1K\_V2}, que chega a 75,3\% de acurácia top-1 no ImageNet com
0,22~GFLOPs por imagem, o menor custo entre as cinco redes do trabalho. A
cabeça original termina numa camada linear de 1280 para 1000 classes
(\texttt{classifier[3]}), que trocamos por uma de 1280 para 5.

Com a troca, o modelo fica com 4.208.437 parâmetros. Com a base congelada, só
os 6.405 parâmetros da cabeça são treinados (0,15\% do total). Por ser leve, é
a rede mais indicada para rodar na câmera do projeto SALUS.
